\clearpage\begingroup
\SourceLeaf{216}\NoSourceHead
\begin{center}
[II.]\par\medskip
\large PRAECEPTA UNIVERSALIA AD COMPUTANDUM DIEM PASCHATIS\par
ANNI CUIUSLIBET DATI SECUNDUM CALENDARIUM TUM\par
GREGORIANUM TUM IULIANUM.
\end{center}
\noindent\rule{\textwidth}{0.4pt}
\begin{center}\small
[Handschriftliche Eintragung in \textsc{Christian Wolf}, Elementa matheseos universae, tomus IV. --- Von \textsc{Gauss} 1800 erworben.]
\end{center}
\noindent\rule{\textwidth}{0.4pt}
\begin{center}
\setlength{\tabcolsep}{5pt}\renewcommand{\arraystretch}{1.3}
\begin{tabular}{|c|c|c|}\hline
Prodeat ex divisione & per & residuum (minimum positivum)\\
Numeri anni & $19$ & $a$\\
Numeri anni & $4$ & $b$\\
Numeri anni & $7$ & $c$\\
$19a+M$ & $30$ & $d$\\
$2b+4c+6d+N$ & $7$ & $e$\\\hline
\multicolumn{3}{|c|}{Eritque dies Paschatis $22+d+e$ Mart. sive $d+e-9$ April.}\\\hline
\end{tabular}
\end{center}
\begin{center}
Ex[emplum] pro $1715$.
\[
\begin{array}{rc}
a&5\\
b&3\\
c&0\\
19a+M&118\\
d&28\\
2b+4c+6d+N&177\\
e&2
\end{array}
\]
Pascha $21$ Apr.
\end{center}
\noindent\hspace*{\parindent}In Calendario Iuliano semper est $M=15$, $N=6$.
\begin{flushright}$27^{\ast}$\end{flushright}

\clearpage\SourceLeaf{217}\HistoricalHead{212}{L}{CHRONOLOGIE.\quad NACHLASS.}
In Gregoriano vero hi numeri per centenos saltem annos valores invariabiles habent; et quidem est
\begin{center}\small
\setlength{\tabcolsep}{3pt}\renewcommand{\arraystretch}{1.25}
\begin{tabular}{rcc|rcc}
Ab Introd. C. G. usque ad $1699$ & $M=22$ & $N=2$ & $2300$---$2399$ & $M=26$ & $N=1$\\
$1700$---$1799$ & $23$ & $3$ & $2400$---$2499$ & $25$ & $1$\\
$1800$---$1899$ & $23$ & $4$ & $2500$---$2599$ & $26$ & $2$\\
$1900$---$1999$ & $24$ & $5$ & $2600$---$2699$ & $27$ & $3$\\
$2000$---$2099$ & $24$ & $5$ & $2700$---$2799$ & $27$ & $4$\\
$2100$---$2199$ & $24$ & $6$ & $2800$---$2899$ & $27$ & $4$\\
$2200$---$2299$ & $25$ & $0$ & $2900$---$2999$ & $28$ & $5$
\end{tabular}\par
etc.
\end{center}
\noindent Generaliter in Calendario Greg[oriano] inveniuntur valores ipsorum $M$, $N$ per centenis annis a $100k$ usque ad $100k+99$ sequenti modo:
\begin{center}
Dividatur $k$ per $\left|\begin{smallmatrix}3\\4\end{smallmatrix}\right|$ et oriatur inde quotiens abiecto residuo $\left|\begin{smallmatrix}p\\q\end{smallmatrix}\right|$ eritque
\end{center}
\[
\begin{aligned}
M&\equiv15+k-p-q\quad(\text{M[od.]}\ 30)\\
N&\equiv4+k-q\quad(\text{M[od.]}\ 7).
\end{aligned}
\]

\textls[120]{Duae ab his regulis dantur Exceptiones in Calend. Greg.}

1) Quoties calculus producit $26^{\mathrm{tum}}$ April. pro die Paschatis, sumi debet ipsius loco $19^{\mathrm{mus}}$ Apr. ---

Hic casus occurrere nequit, nisi in saeculis, ubi $11(M+1)$ per $30$ divisus praebet residuum minus quam $19$, puta ubi $M\equiv0,\allowbreak 2,\allowbreak 3,\allowbreak 5,\allowbreak 6,\allowbreak 8,\allowbreak 10,\allowbreak 11,\allowbreak 13,\allowbreak 14,\allowbreak 16,\allowbreak 17,\allowbreak 19,\allowbreak 21,\allowbreak 22,\allowbreak 24,\allowbreak 25,\allowbreak 27,\allowbreak 29$.

2) Quoties evadit $d=28$, $e=6$ atque insuper $11(M+1)$ per $30$ divisus producit residuum minus quam $8$, Pascha celebratur non $25^{\mathrm{to}}$ April., sed $18^{\mathrm{vo}}$ April. Ad hunc finem esse debet $M=2,\allowbreak 5,\allowbreak 10,\allowbreak 13,\allowbreak 16,\allowbreak 21,\allowbreak 24,\allowbreak 29$.

\begin{center}\small BEMERKUNGEN.\end{center}
\begingroup\small
\noindent\hspace*{\parindent}Die vorliegende Eintragung scheint \textsc{Gauss} erst nach seiner Veröffentlichung \textit{Berechnung des Osterfestes} in der Monatlichen Correspondenz zur Beförderung der Erd- und Himmels-Kunde 2, August 1800, S. 121, Werke VI, S. 73, gemacht zu haben. In der hier vorliegenden Aufzeichnung wird bei dem zweiten Ausnahmefalle bemerkt, daß für die in Frage kommenden $8$ Werte des $M$ der Ausdruck $11(M+1)$ durch $30$ dividiert einen Rest gibt, der sogar kleiner als $8$ ist, nicht nur, wie es in \textsc{Gauss}' Veröffentlichung heißt,
\clearpage\SourceLeaf{218}\HistoricalHead{213}{R}{BEMERKUNGEN.}
\noindent kleiner als $19$ (siehe Werke VI, S. 79, Z. 12 v. u.); ferner werden hier die Werte von $M$ und $N$ bis für das Jahr $2999$ statt wie dort bis $2499$ angegeben.

In \textsc{Kästners} \textit{Anfangsgründen der angewandten Mathematik} (Der mathematischen Anfangsgründe zweyter Theil), Zweyte Abtheilung, von \textsc{Gauss} 1791 erworben, hat \textsc{Gauss} auch eine Eintragung über die Bestimmung des Osterfestes gemacht, die aus der Zeit vor der ersten Veröffentlichung zu stammen scheint, wie aus der noch nicht ganz klaren Angabe über die zwei Ausnahmefälle im Gregorianischen Kalender hervorgeht. Hingegen ist dort eine andere Bestimmung von $M$ und $N$ angegeben. In seiner ersten Veröffentlichung (Werke VI, S. 78) gibt \textsc{Gauss}, ebenso wie oben S. 212, zur Bestimmung von $M$ und $N$ folgende Vorschrift:
\begin{center}
$\left\{\begin{smallmatrix}M\\N\end{smallmatrix}\right\}$ der Rest, den man erhält, wenn man $\left\{\begin{smallmatrix}15+k-p-q\\4+k-q\end{smallmatrix}\right\}$ mit $\left\{\begin{smallmatrix}30\\7\end{smallmatrix}\right\}$ dividiert,
\end{center}
\noindent dabei ist
\[
p=Q\left(\frac{k}{3}\right),\quad q=Q\left(\frac{k}{4}\right)
\]
\noindent zu nehmen. In der erwähnten Eintragung bei \textsc{Kästner} heißt es:
\begin{quote}\normalsize
»Allgemein von $100k$ bis $100k+99$
\[
\begin{aligned}
k&\equiv p\quad(\text{Mod.}\ 3)\\
&\equiv q\quad(\text{Mod.}\ 4)\\
&\equiv r\quad(\text{Mod.}\ 7)\\
M&\equiv15+\tfrac{5}{12}k+\tfrac{1}{3}p+\tfrac{1}{4}q\quad(\text{Mod.}\ 30)\\
N&\equiv4+2q+6r\quad(\text{Mod.}\ 7)\text{«.}
\end{aligned}
\]
\end{quote}
\noindent\hspace*{\parindent}Bedeuten $Q$ und $R$ Quotient bezw. Rest und bedient man sich der Identitäten
\begingroup\makeatletter\tagsleft@true\let\veqno\@@leqno\makeatother
\begin{align}
k&=3Q\left(\frac{k}{3}\right)+R\left(\frac{k}{3}\right),\tag*{(1)}\\
k&=4Q\left(\frac{k}{4}\right)+R\left(\frac{k}{4}\right),\tag*{(2)}\\
k&=7Q\left(\frac{k}{7}\right)+R\left(\frac{k}{7}\right),\tag*{(3)}
\end{align}
\endgroup
\noindent so ist erstens nach (1) und (2)
\[
\begin{aligned}
15+k-Q\left(\frac{k}{3}\right)-Q\left(\frac{k}{4}\right)
&=15+k-\frac{k-R\left(\frac{k}{3}\right)}{3}-\frac{k-R\left(\frac{k}{4}\right)}{4}\\
&=15+\frac{5}{12}k+\frac{1}{3}R\left(\frac{k}{3}\right)+\frac{1}{4}R\left(\frac{k}{4}\right);
\end{aligned}
\]
\noindent multipliziert man ferner (2) mit $2$ und (3) mit $6$ und addiert, so hat man
\[
8k=8Q\left(\frac{k}{4}\right)+2R\left(\frac{k}{4}\right)+42Q\left(\frac{k}{7}\right)+6R\left(\frac{k}{7}\right)
\]
\noindent oder
\[
k\equiv Q\left(\frac{k}{4}\right)+2R\left(\frac{k}{4}\right)+6R\left(\frac{k}{7}\right)\quad(\text{mod.}\ 7).
\]
\noindent Folglich ist zweitens
\[
4+k-Q\left(\frac{k}{4}\right)\equiv4+2R\left(\frac{k}{4}\right)+6R\left(\frac{k}{7}\right)\quad(\text{mod.}\ 7).
\]
\clearpage\SourceLeaf{219}\HistoricalHead{214}{L}{CHRONOLOGIE.\quad NACHLASS.}
\noindent Hiermit sind beide Aussagen von \textsc{Gauss} als gleichwertig erwiesen; beide haben nur bis zum Jahre $4200$ Gültigkeit (vergl. S. 201 oben).

Zu \textsc{Gauss}' erster Veröffentlichung über das Osterfest (Werke VI, S. 73) ist noch zu bemerken, daß er auf S. 128--129 seines Exemplars der Monatlichen Correspondenz zur Beförderung der Erd- und Himmels-Kunde 2, zwischen dem Text seinem Aufsatze handschriftliche Bemerkungen beigefügt hat. Von diesen hat \textsc{Schering} bei dem Abdruck des \textsc{Gauss}schen Aufsatzes in die Werke VI, S. 79, Z. 11 die folgende übernommen »(Z. B. 1609, 1989)«. Bei der letzten Zahl hat sich \textsc{Gauss} verschrieben; die fragliche Jahreszahl ist 1981, wie man Werke VI, S. 85, Z. 15 v. u. richtig angegeben findet. Auch in der erwähnten Eintragung in \textsc{Kästners} \textit{Anfangsgründen} heißt es in korrekter Weise bei den Ausnahmefällen:
\begin{center}\normalsize
Dies geschieht zum ersten Male $\begin{array}{r}1954\\1981,\end{array}$ wo Ostern den $\begin{array}{c}18_{\mathrm{ten}}\\19\end{array}$ April fällt.
\end{center}
\noindent\hspace*{\parindent}Ferner hat \textsc{Gauss} in sein Exemplar der Monatlichen Correspondenz bei der Ausnahme II (Werke VI, S. 79, Z. 18) eingetragen:
\begin{center}\normalsize
oder ist $a>10$.
\end{center}
\noindent Dieses einfachere Kriterium statt des Restcharakters von $11M+11$ (vergl. oben S. 212) findet sich auch in der \textsc{Gauss}schen Veröffentlichung \textit{Noch etwas über die Bestimmung des Osterfestes} (Werke VI, S. 85, Z. 12 v. u.). Die weiteren \textsc{Gauss}schen Eintragungen am besprochenen Orte hat \textsc{Schering} unter der Überschrift »[Handschriftliche Bemerkung]« in die Werke VI, S. 79 unten übernommen; vergl. zu der dort gegebenen Bestimmung von $p$ die \textit{Berichtigung}, oben S. 201, zu den auf die Zeiträume zwischen Ostern und Michaelis bezüglichen Zählungen \textsc{Franz Goldscheider}, \textit{Über die Gausssche Osterformel} II, Wissenschaftliche Beilage zum Jahresbericht des Luisenstädtischen Realgymnasiums zu Berlin, Ostern 1899, S. 29.
\begin{flushright}\textsc{Loewy.}\end{flushright}
\endgroup
\begin{center}\rule{0.25\textwidth}{0.4pt}\end{center}
\clearpage\endgroup
